\chapter{数値計算の方法と検証}\label{app:num}

\section{数値計算の方法}

\begin{longtable}{p{0.16\linewidth}p{0.77\linewidth}}
\toprule
対象 & 方法 \\ \midrule \endhead
運動学（$L=0$） & 式~\eqref{eq:basic} を常微分方程式として積分。Sim01・02 は適応刻みの Runge--Kutta 法（相対誤差 $10^{-10}$）、Sim03 は 4 次 Runge--Kutta 法（刻み 2 ms）、Sim05A は中点法（刻み 40 ms、接続を格子上に事前計算して双一次補間）。 \\
接続 $\Gamma$ と曲率 $F$ & 角運動量が速度に線形であることを使い、単位速度を入れて $A$, $B$ を求める。$F$ は格子上の中心差分、または中心差分で各点を計算。 \\
順動力学（Sim04） & 重心系の質量行列 $M(q)$ を各剛体のヤコビアンから組み立て、$\partial M/\partial q$ は中心差分（刻み $10^{-6}$）。積分は適応刻みの Runge--Kutta 法。 \\
3 次元（Sim07・09） & 腕・脚を 40 個の点質量で表し、$\bm\omega=-A^{-1}B\dot q$ を中点則で積分、姿勢は回転ベクトルの指数写像で更新。 \\
最適化 & 軌道最適化は SLSQP \cite{R26}（多点から開始し、制約を満たす最良解を採用）。推進剤配分と段差の摩擦は線形計画（HiGHS \cite{R25}）。 \\
学習（Sim05） & OpenAI-ES \cite{R19}：対称サンプリング、順位変換した適応度、Adam \cite{R20} による更新。方策は tanh を活性化関数とする全結合ネットワーク。母集団全体を行列演算で同時に評価。 \\
地上の反力（Sim06・07・09） & 関節経路を時間 $\tau\in[0,1]$ の等間隔点で与え、位置を数値的に 2 回微分して加速度を求め、式~\eqref{eq:zmp} で ZMP を計算。動作時間 $T$ に対しては加速度が $1/T^2$ で変わることを使う。 \\
スティックスリップ（Sim06） & 静止摩擦・動摩擦を区別する Karnopp 型のモデル \cite{R30}、刻み 0.2 ms の陽解法。 \\
\bottomrule
\end{longtable}

\section{検証}

モデルと数値計算の正しさを、独立した方法どうしの一致で確認した。

\begin{longtable}{p{0.40\linewidth}p{0.53\linewidth}}
\toprule
検証の内容 & 結果 \\ \midrule \endhead
腕を往復させた後の胴体の姿勢（Sim01） & 残留 $5\times10^{-9}$°（可逆性） \\
人間大と 18 m 級（長さ 10 倍・質量 1000 倍）の回転角（Sim01） & 両方とも $-16.0432$°（スケール不変性） \\
ループ 1 周の回転：数値積分とストークスの定理（Sim02） & $+2.9778$° で一致 \\
トルクだけを入力した順動力学と運動学（Sim04） & 2 周後の回転 $+13.4086$° で一致、関節軌道の誤差 $2.6\times10^{-8}$° \\
順動力学の角運動量・エネルギー収支（Sim04） & $|L|\le 2\times10^{-9}$ kg\,m$^2$/s、$|\Delta T - \int P\,dt| \le 9\times10^{-4}$ J \\
スケーリング則の指数（Sim04） & トルク $s^{5.00}$、$f_{\max}$ $s^{-1.00}$、容量 $s^{4.00}$（理論どおり） \\
平均化モデルと位相モデル（Sim05B） & 60 s 後の周回平均の姿勢 89.63° 対 90.00° \\
格子補間と厳密な接続（Sim05A） & 最大差 0.32°（学習方策の軌道を再生して確認） \\
フルード則（Sim06） & $T_{\min}/\sqrt{s\,g_{\text{地球}}/g}$ が 12 条件で 1.616 \\
3 次元モデルと平面モデル（Sim07） & 腕を平面内で動かした場合に 5 桁一致 \\
円錐運動の回転の指数（Sim07） & 半頂角の 2.00 乗（2 次効果） \\
コッドマンのパラドックス（Sim08） & ひねりと立体角が全例で小数第 2 位まで一致 \\
\bottomrule
\end{longtable}

\section{修正の記録}

\begin{itemize}
\item \textbf{Sim06・Sim09 の関節制約}：腕の軌道最適化で関節速度・加速度の制約を間引いて課す際、配列の取り方の誤りにより、肘の関節に制約がかかっていなかった。
  修正後、Sim06 の地球での最短時間は 1.00 s（38\% 短縮）から 1.03 s（36\% 短縮）に変わった。月面の結果と Sim09 の結論は変わらない。
\item \textbf{Sim05B の揺れ補償 MPC}：当初は目標を「現在の位相での揺れ」で補正していたが、必要なのは「旋回終了時の位相での揺れ」であった。計画から終了時の位相を予測するよう改めた。
\item \textbf{Sim08 の腕組みの判定}：人の可動域の内旋 70° は腕を 90° 外転した姿勢での標準値であり、ひねりの値は測り方で変わる（第~\ref{ch:shoulder}章）。
  そのため、合否の判定に加えて「最低限必要な内旋量」を求める分析を行った。
\item \textbf{Sim09 の比較の公平性}：二本脚にだけ上体の傾きを許していたため、戦車型にも腰の関節を与え、上体の速度・加速度を腕より厳しく制限した。
\item \textbf{Sim10 の二本脚}：当初は膝だけで体を持ち上げる悲観的なモデルだったため、上体の前傾による重心移動と股関節トルクを加えた。
\end{itemize}

\section{結果の限界}

\begin{itemize}
\item 平面モデルは人の体を矢状面に投影した粗いモデルで、左右の手足の区別や体のねじれは扱っていない。
\item 関節トルクの上限・可動域は目安の値であり、個人差・機体差は扱っていない。
\item 地上の評価は準静的（台車が静止した状態での転倒・滑りの開始条件）で、転倒が始まった後の運動は扱っていない。
\item 学習の結果は乱数の種と超パラメータに依存しうる。主な結論（構造の有無による差）は再実行でも同じであった。
\end{itemize}
